\documentclass[11pt,a4paper]{article}
\usepackage{amsmath,amssymb}
\usepackage{listings}
\usepackage{hyperref}
\usepackage[margin=1in]{geometry}
\usepackage{fancyhdr}

\pagestyle{fancy}
\fancyhf{}
\fancyhead[L]{Module 09: File I/O}
\fancyhead[R]{C Programming Course}
\fancyfoot[C]{\thepage}

\lstset{
  language=C,
  basicstyle=\ttfamily\small,
  keywordstyle=\bfseries,
  commentstyle=\itshape,
  numbers=left,
  numberstyle=\tiny,
  frame=single,
  breaklines=true,
  tabsize=4
}

\title{Module 09: File I/O}
\author{C Programming Course}
\date{}

\begin{document}
\maketitle
\thispagestyle{fancy}

\section{File Handling Basics}
Files are accessed through \texttt{FILE *} pointers using \texttt{fopen} and
\texttt{fclose}. Common modes: \texttt{"r"} (read), \texttt{"w"} (write),
\texttt{"a"} (append), \texttt{"rb"}/\texttt{"wb"} (binary).

\section{Reading and Writing Text Files}
\begin{lstlisting}
#include <stdio.h>

int main(void) {
    FILE *fp = fopen("output.txt", "w");
    if (fp == NULL) {
        perror("fopen");
        return 1;
    }
    fprintf(fp, "Hello, File I/O!\n");
    fprintf(fp, "Line number: %d\n", 2);
    fclose(fp);

    /* Read the file back */
    fp = fopen("output.txt", "r");
    if (fp == NULL) {
        perror("fopen");
        return 1;
    }
    char buf[128];
    while (fgets(buf, sizeof(buf), fp) != NULL) {
        printf("%s", buf);
    }
    fclose(fp);
    return 0;
}
\end{lstlisting}

\section{Binary File Operations}
\begin{lstlisting}
#include <stdio.h>

int main(void) {
    int data[4] = {10, 20, 30, 40};
    FILE *fp = fopen("data.bin", "wb");
    if (fp == NULL) return 1;
    fwrite(data, sizeof(int), 4, fp);
    fclose(fp);

    int result[4];
    fp = fopen("data.bin", "rb");
    if (fp == NULL) return 1;
    fread(result, sizeof(int), 4, fp);
    fclose(fp);
    return 0;
}
\end{lstlisting}

\section{File Positioning}
\begin{itemize}
  \item \texttt{fseek(fp, offset, whence)} -- moves the file position.
  \item \texttt{ftell(fp)} -- returns the current position.
  \item \texttt{rewind(fp)} -- resets position to the beginning.
\end{itemize}

Whence values: \texttt{SEEK\_SET} (beginning), \texttt{SEEK\_CUR} (current),
\texttt{SEEK\_END} (end).

\section{Error Handling in File Operations}
\begin{itemize}
  \item Always check if \texttt{fopen} returns \texttt{NULL}.
  \item Use \texttt{perror} or \texttt{strerror(errno)} for diagnostics.
  \item Check return values of \texttt{fread}/\texttt{fwrite} for partial I/O.
  \item Use \texttt{ferror(fp)} to detect stream errors.
\end{itemize}

\end{document}
